% from ega0-16-fr.tex
\oldpage[0\textsubscript{IV-1}]{33}
maximal (Bourbaki, \emph{Alg.~comm.}, chap.~II, \S~1, no~1, prop.~2)~;
autrement dit, il y a un élément $z\in N$ dont l'annulateur est
$\mathfrak{m}$, et par suite le sous-module $Az$ de $N$ est isomorphe à
$A/\mathfrak{m}=k$~; d'où le lemme.
\end{proof}

\begin{lemma}[16.4.3]
\label{0.16.4.3-fr}
Soient $A$ un anneau local noethérien, $\mathfrak{m}$ son idéal maximal,
$k=A/\mathfrak{m}$ son corps résiduel, $M$ un $A$-module,
$(x_i)_{1\leq i\leq n}$ une suite $M$-régulière d'éléments de
$\mathfrak{m}$. Alors les $A$-modules
\[
  \operatorname{Hom}_A(k,M/(x_1M+\cdots+x_nM))
  \quad\text{et}\quad \operatorname{Ext}^n_A(k,M)
\]
sont isomorphes, et pour $n\geq1$, ils sont aussi isomorphes à
$\operatorname{Ext}^{n-1}_A(k,M/x_1M)$.
\end{lemma}

\begin{proof}
Raisonnons par récurrence sur $n$, la proposition étant évidente pour $n=0$.
Posons $N=M/x_1M$~; la suite $(x_i)_{2\leq i\leq n}$ est $N$-régulière, donc
\[
  \operatorname{Hom}_A(k,M/(x_1M+\cdots+x_nM))
  =\operatorname{Hom}_A(k,N/(x_2N+\cdots+x_nN))
\]
est isomorphe à $\operatorname{Ext}^{n-1}_A(k,N)$ en vertu de l'hypothèse de
récurrence. Considérons alors la suite exacte
$0\to M\xrightarrow{x_1}M\to M/x_1M=N\to0$~; on en déduit la suite exacte des
Ext
\begin{equation}
\label{0.16.4.3.1-fr}
\tag{16.4.3.1}
  \operatorname{Ext}^{n-1}_A(k,M)
  \to\operatorname{Ext}^{n-1}_A(k,N)
  \to\operatorname{Ext}^n_A(k,M)
  \xrightarrow{x_1}\operatorname{Ext}^n_A(k,M).
\end{equation}
En vertu de l'hypothèse de récurrence, $\operatorname{Ext}^{n-1}_A(k,M)$ est
isomorphe à
\[
  \operatorname{Hom}_A(k,M/(x_1M+\cdots+x_{n-1}M)),
\]
qui est nul, puisque $(x_i)_{1\leq i\leq n-1}$ n'est pas une suite
$M$-régulière maximale (16.4.2). D'autre part, comme
$x_1\in\mathfrak{m}$, l'homothétie de rapport $x_1$ dans le $A$-module
$\operatorname{Hom}_A(k,T)$ est nulle pour tout $A$-module $T$, donc il en est
de même de l'homothétie de rapport $x_1$ dans tout $A$-module
$\operatorname{Ext}^n_A(k,T)$~; les assertions du lemme découlent donc
aussitôt de la suite exacte (16.4.3.1).
\end{proof}

\begin{corollary}[16.4.4]
\label{0.16.4.4-fr}
Soient $A$ un anneau local noethérien, $\mathfrak{m}$ son idéal maximal,
$k=A/\mathfrak{m}$ son corps résiduel, $M$ un $A$-module de type fini
$\neq0$. Toutes les suites $M$-régulières maximales d'éléments de
$\mathfrak{m}$ ont le même nombre d'éléments, qui est le plus petit entier $n$
tel que $\operatorname{Ext}^n_A(k,M)\neq0$.
\end{corollary}

\begin{definition}[16.4.5]
\label{0.16.4.5-fr}
Soient $A$ un anneau local noethérien, $\mathfrak{m}$ son idéal maximal, $M$ un
$A$-module de type fini. On appelle \emph{profondeur} de $M$, et l'on note
$\operatorname{prof}_A M$ ou $\operatorname{prof}(M)$ la borne supérieure du
nombre d'éléments d'une suite $M$-régulière formée d'éléments de
$\mathfrak{m}$. On a donc $\operatorname{prof}(M)=+\infty$ si et seulement si
$M=0$, et, par (16.4.1)
\begin{equation}
\label{0.16.4.5.1-fr}
\tag{16.4.5.1}
  \operatorname{prof}(M)\leq\dim(M)\quad\text{si}\quad M\neq0.
\end{equation}
On a évidemment $\operatorname{prof}(M^k)=\operatorname{prof}(M)$ pour tout
$k>0$.
\end{definition}

\begin{proposition}[16.4.6]
\EGAReaderTarget{canon-ref.0.16.4.6-fr}{0.16.4.6-fr}
Soient $A$ un anneau local noethérien, $\mathfrak{m}$ son idéal maximal, $M$ un
$A$-module de type fini.
\begin{enumerate}
  \item[\rm{(i)}] Pour que $\operatorname{prof}(M)=0$, il faut et il suffit que
  $\mathfrak{m}\in\operatorname{Ass}(M)$ (ou encore qu'il existe un sous-module
  de $M$ isomorphe au corps résiduel $k=A/\mathfrak{m}$ de $A$).
  \item[\rm{(ii)}] Pour tout élément $M$-régulier $x$ appartenant à
  $\mathfrak{m}$, on a
  \begin{equation}
  \label{0.16.4.6.1-fr}
  \tag{16.4.6.1}
    \operatorname{prof}(M/xM)=\operatorname{prof}(M)-1.
  \end{equation}
  \item[\rm{(iii)}] Si $M\neq0$, on a
  \begin{equation}
  \label{0.16.4.6.2-fr}
  \tag{16.4.6.2}
    \operatorname{prof}(M)
    \leq\inf_{\mathfrak{p}\in\operatorname{Ass}(M)}\dim(A/\mathfrak{p})
    \leq\dim(M).
  \end{equation}
  